\chapter{The exact unbounded-state-scale C2 complement}
\label{component:pb60046a37191}
\noindent\textbf{Component }\path{pb60046a37191}\par
{\small\noindent Source: \path{R13/baseline/baseline/baseline/baseline/advances/source/unbounded_b_reduction.tex}\par}
{\small\noindent SHA256: \path{b60046a37191c59da6ebca1e87e245f5c8df75d7c802327c1cb821c8629f81a3}\par}
{\small\noindent Inventory use: SR1 exact opening-cancelled source scales and exhaustive B=xb weak/unit split\par}
\medskip
\section[The exact unbounded-state-scale C2 complement]{The exact unbounded-state-scale C2 complement}
\label{pb60046a37191:sec:c2-unbounded-b-reduction}

Use the original positive rational source with $d=1$, actual
$u_c=\sqrt x$ positive infinitesimal and $v_c=k$ a positive bounded
unit. The global entry theorem supplies its positive chart inequalities.
The preceding rational theorems, including the joint-opening theorem,
now cover every bounded $b$: the inherited $n\ge0$ entry inequality
makes bounded $b$ a unit. We classify the remaining $b$ unbounded
without changing any original coefficient.

Put
\[
 B=xb,\qquad H=xh,\qquad L=x(b-h)=B-H,\qquad
 a^2=k^2/x,\quad t=L/k^2.
\]
The original gaps give
\[
 0<H<1,\qquad0<L<k^2(1+xr),\qquad B=H+L.
\]
Thus $B,H,L$ are bounded and $x/B=1/b$ is infinitesimal.
Set $A_c=1+xr$ and $S=A_c+L$, both positive bounded units. Also
\[
 g/b=1+(r-1)/b+H/b\longrightarrow1,
 \qquad j=1+xr-L/k^2>0,\quad j\le1+x.
\]
The following exact rearrangements of the original short polynomials
control every event scale:
\begin{align*}
 x^2N={}&(1+x)(B+xr)(H+k^2)-x(1+k^2)(1-r)L,\\
 \Delta/b={}&(1+x)B+2x(1+x)r+(1+x)(1+k^2)\\
 &+(x^2-k^2x)r^2/b+2x(1+k^2)r/b+(1+k^2)/b.
\end{align*}
Consequently $xN/b=x^2N/B$ and $\Delta/b$ are positive bounded
units, with respective residues
\[
 H_0+k_0^2,\qquad B_0+1+k_0^2.
\]
Here $H,L\le B$ controls the divided corrections, and $x/B\to0$.

The exact physical condition $n\ge0$ is
\[
 r(\Delta/b)\le(1+r/b)A_cS.
\]
It implies
\[
 r_0\le\frac{1+L_0}{1+k_0^2+B_0},\qquad
 1-r_0\ge\frac{k_0^2+H_0}{1+k_0^2+B_0}>0.
\]
Thus $1-r$ is a positive bounded unit without any lower bound on $r$.
The original formulas now give the exact comparabilities
\[
 m\asymp b^{-2},\qquad C\asymp L^{-1},\qquad A\asymp H^{-1},
 \qquad D\asymp B/j,\qquad\alpha\asymp B/(1-H).
 \tag{SR1}
\]
Every omitted factor in $\asymp$ is a positive bounded unit with
bounded reciprocal. For instance,
\[
 CL=\frac{A_c(xN/b)}{2\Delta/b},\qquad
 AH=\frac{k^2A_c(1+r/b)S}{2\Delta/b}
\]
are positive units; the other identities follow using $g/b\to1$
and $(b-k^2r)/b\to1$.

There are exactly two unbounded-$b$ scale leaves:
\begin{enumerate}
\item If $B=xb$ is infinitesimal, so are $H,L$, while $j\to1$
and $1-H\to1$. Thus $C,A$ are unbounded and $m,D,\alpha$ are
weak, with $m\asymp b^{-2}$ and $D\asymp\alpha\asymp xb$.
This is the one-weak-complementary-event frontier, with weak
$\alpha$ and strong $C$ in this actual-angle orientation.
\item If $B$ is a positive bounded unit, then $C,A,D,\alpha$
are all strong and $m$ is weak. The new strong-complement pitchfork
applies to this entire leaf at all upper event scales. In particular
it includes $C$ unbounded when $L$ is infinitesimal and $\alpha$
unbounded when $1-H$ is infinitesimal.
\end{enumerate}
Unbounded $B$ is impossible by the original chart inequalities.
The homogeneous all-tests theorem of Section~\ref{pab6222b899ef:sec:c2-unbounded-b-small-B}
closes the remaining infinitesimal-$B$ leaf without assuming a bounded
actual state ratio. Thus both unbounded-$b$ leaves, as well as all
bounded-$b$ openings, are now covered. Literal $n=0$ is preserved.
